% =====================================================================
% 07_partie_e.tex — fonctionnalités objet avancées
% =====================================================================
\chapter{Partie E : fonctionnalités objet avancées}
\label{chap:partieE}

\section{Objectifs de la partie E}

La partie E regroupe les mécanismes objet les plus avancés du projet : un type
récursif, une collection multi-niveaux et une vue objet avec méthode. Elle
prolonge la partie A, qui a préparé les types, et la partie D, qui a fourni les
méthodes. Elle ne crée pas de nouvelle table métier : elle exploite les
structures existantes pour montrer la navigation objet, l'agrégation de
collections et la mise à jour d'affectations par une méthode de vue.

\begin{objectifbox}[Trois démonstrations]
E10 montre le parrainage récursif et la double notation
\code{c.ref\_parrain.ref\_parrain.nom}. E11 montre un \code{VARRAY} de mois
contenant une \emph{nested table} de statistiques. E12 montre une vue objet
dont les collections sont calculées par \code{CAST(MULTISET(...))} et dont la
méthode \code{mutation} modifie l'historique des affectations.
\end{objectifbox}

\section{E10 : type récursif et parrainage}

L'attribut \code{ref\_parrain} a été ajouté à \code{client\_type} en partie A
par \code{ALTER TYPE ... ADD ATTRIBUTE ... CASCADE}. Il est de type
\code{REF client\_type} : un client peut donc référencer un autre client, et la
chaîne peut être parcourue sur plusieurs niveaux. La partie A a déjà parrainé
les clients 2 et 3 par le client 1. E10 ajoute trois liens supplémentaires :
client 4 vers 1, client 5 vers 2 et client 6 vers 3. Les clients 5 et 6
obtiennent ainsi un grand-parrain.

\begin{lstlisting}[caption={E10 : mise à jour des liens de parrainage.},label={lst:e10}]
UPDATE client
   SET ref_parrain = (SELECT REF(p) FROM client p WHERE p.num_client = 1)
 WHERE num_client = 4;
UPDATE client
   SET ref_parrain = (SELECT REF(p) FROM client p WHERE p.num_client = 2)
 WHERE num_client = 5;
UPDATE client
   SET ref_parrain = (SELECT REF(p) FROM client p WHERE p.num_client = 3)
 WHERE num_client = 6;
COMMIT;
\end{lstlisting}

La requête de contrôle utilise deux fois la notation pointée. Le premier
\code{ref\_parrain} atteint le parrain direct ; le second atteint son propre
parrain. Oracle suit la chaîne parce que \code{client\_type} contient
lui-même l'attribut \code{ref\_parrain}. Un alias de table est indispensable :
sans \code{c}, Oracle ne sait pas à quelle ligne appliquer le déréférencement
et lève \code{ORA-00904}.

\begin{lstlisting}[caption={E10 : lecture du parrain et du grand-parrain.},label={lst:e10-lecture}]
SELECT c.num_client, c.nom,
       c.ref_parrain.nom             AS nom_parrain,
       c.ref_parrain.ref_parrain.nom AS nom_grand_parrain
  FROM client c
 WHERE c.ref_parrain IS NOT NULL
 ORDER BY c.num_client;
\end{lstlisting}

\begin{table}[htbp]
  \centering
  \small
  \caption{Résultats du parrainage après E10.}
  \label{tab:resultats-e10}
  \begin{tabularx}{\textwidth}{@{}C{1.5cm}L{4.0cm}L{4.0cm}X@{}}
    \toprule
    \entete{Client} & \entete{Parrain} & \entete{Grand-parrain} &
    \entete{Commentaire} \\
    \midrule
    2 & AMRANI & --- & Parrain direct. \\
    3 & AMRANI & --- & Parrain direct. \\
    4 & AMRANI & --- & Parrain direct. \\
    5 & BENALI & AMRANI & Chaîne sur deux niveaux. \\
    6 & CHERIF & AMRANI & Chaîne sur deux niveaux. \\
    \bottomrule
  \end{tabularx}
\end{table}

\begin{clebox}[Pourquoi le type récursif est-il important ?]
Un type récursif permet de représenter une relation hiérarchique sans table
d'association supplémentaire. Ici, le parrainage est une relation orientée :
chaque client possède au plus un parrain, mais peut avoir plusieurs filleuls.
La double notation montre que la référence est réellement navigable, et non un
simple identifiant stocké sous forme de texte.
\end{clebox}

\section{E11 : collection multi-niveaux \code{releves}}

La table \code{releves} contient une année et un \code{VARRAY(12)} de mois.
Chaque mois contient un numéro, un nombre d'opérations et une \emph{nested
table} de statistiques par type d'opération. La structure est donc
hiérarchique sur trois niveaux : année, mois, type d'opération.

\begin{table}[htbp]
  \centering
  \small
  \caption{Structure de la collection \code{releves}.}
  \label{tab:structure-releves}
  \begin{tabularx}{\textwidth}{@{}L{4.5cm}L{4.6cm}X@{}}
    \toprule
    \entete{Élément} & \entete{Type} & \entete{Contenu} \\
    \midrule
    \code{releves} & Table relationnelle &
    \code{annee NUMBER PRIMARY KEY}, \code{mois\_vry mois\_vry\_type}. \\
    \code{mois\_vry\_type} & \code{VARRAY(12)} &
    Douze mois au maximum. \\
    \code{mois\_elt\_vry\_type} & Objet &
    \code{mois}, \code{nb\_operations}, \code{types}. \\
    \code{type\_stat\_nt\_type} & \emph{Nested table} &
    Collection de \code{type\_stat\_elt\_nt\_type}. \\
    \code{type\_stat\_elt\_nt\_type} & Objet &
    \code{type\_op}, \code{nb}, \code{montant\_total}. \\
    \bottomrule
  \end{tabularx}
\end{table}

Quatre mois sont d'abord insérés, puis un cinquième est ajouté par
\code{EXTEND(1)}. La nouvelle \emph{nested table} est construite avec un
constructeur \code{type\_stat\_nt\_type(...)} et affectée au dernier élément
du \code{VARRAY}. Le bloc affiche le nombre de mois avant et après extension,
ainsi que le nombre de types du mois ajouté.

\begin{lstlisting}[caption={E11 : insertion de quatre mois.},label={lst:e11-insert},basicstyle=\ttfamily\scriptsize]
INSERT INTO releves (annee, mois_vry)
VALUES (2026, mois_vry_type(
  mois_elt_vry_type(1, 10, type_stat_nt_type(
    type_stat_elt_nt_type('RETRAIT', 6, 120000),
    type_stat_elt_nt_type('DEPOT',   3, 200000),
    type_stat_elt_nt_type('CB',      1,  15000))),
  mois_elt_vry_type(2, 8, type_stat_nt_type(
    type_stat_elt_nt_type('RETRAIT', 5,  90000),
    type_stat_elt_nt_type('DEPOT',   2, 150000),
    type_stat_elt_nt_type('VIREMENT',1,  50000))),
  mois_elt_vry_type(3, 12, type_stat_nt_type(
    type_stat_elt_nt_type('RETRAIT', 7, 140000),
    type_stat_elt_nt_type('DEPOT',   4, 250000),
    type_stat_elt_nt_type('CB',      1,  20000))),
  mois_elt_vry_type(4, 9, type_stat_nt_type(
    type_stat_elt_nt_type('RETRAIT', 6, 110000),
    type_stat_elt_nt_type('DEPOT',   2, 180000),
    type_stat_elt_nt_type('VIREMENT',1,  30000)))
));
COMMIT;
\end{lstlisting}

\begin{lstlisting}[caption={E11 : extension à un cinquième mois.},label={lst:e11-extend}]
DECLARE
  v_mois  mois_vry_type;
  v_types type_stat_nt_type;
  v_av    NUMBER;
  v_ap    NUMBER;
BEGIN
  SELECT mois_vry INTO v_mois FROM releves WHERE annee = 2026 FOR UPDATE;

  v_av := v_mois.COUNT;
  v_mois.EXTEND(1);
  v_types := type_stat_nt_type(
    type_stat_elt_nt_type('RETRAIT', 8, 160000),
    type_stat_elt_nt_type('DEPOT',   2, 100000),
    type_stat_elt_nt_type('VIREMENT',1,  50000));
  v_mois(v_mois.LAST) := mois_elt_vry_type(5, 11, v_types);

  UPDATE releves SET mois_vry = v_mois WHERE annee = 2026;

  v_types := v_mois(v_mois.LAST).types;
  v_ap := v_mois.COUNT;
  DBMS_OUTPUT.PUT_LINE('VARRAY : '||v_av||' -> '||v_ap||' mois');
  DBMS_OUTPUT.PUT_LINE('Nested table du mois '||v_mois(v_mois.LAST).mois
    ||' : '||v_types.COUNT||' types');
  COMMIT;
END;
/
\end{lstlisting}

L'interrogation nécessite deux appels à \code{TABLE} : le premier ouvre le
\code{VARRAY} des mois, le second ouvre la \emph{nested table} des types. La
syntaxe \code{TABLE(m.COLUMN\_VALUE.types)} ne fonctionne pas ; il faut écrire
\code{TABLE(m.types)}. De même, \code{CARDINALITY} sur \code{mois\_vry} lève
\code{ORA-00932} dans cette version d'Oracle ; on utilise \code{TABLE(...)} et
\code{COUNT(*)}.

\begin{lstlisting}[caption={E11 : total annuel par type d'opération.},label={lst:e11-total}]
SELECT t.type_op,
       SUM(t.nb) AS nb_total,
       SUM(t.montant_total) AS montant_total
  FROM releves r, TABLE(r.mois_vry) m, TABLE(m.types) t
 WHERE r.annee = 2026
 GROUP BY t.type_op
 ORDER BY montant_total DESC;
\end{lstlisting}

\begin{table}[htbp]
  \centering
  \small
  \caption{Totaux du relevé 2026 sur cinq mois.}
  \label{tab:resultats-e11}
  \begin{tabularx}{0.62\textwidth}{@{}L{3.4cm}C{2.6cm}C{3.2cm}@{}}
    \toprule
    \entete{Type} & \entete{Nombre total} & \entete{Montant total (DA)} \\
    \midrule
    DEPOT & 13 & 880\,000 \\
    RETRAIT & 32 & 620\,000 \\
    VIREMENT & 3 & 130\,000 \\
    CB & 2 & 35\,000 \\
    \bottomrule
  \end{tabularx}
\end{table}

\begin{remarquebox}[Lecture des totaux]
Les totaux portent sur les cinq mois, pas sur les opérations de la table
\code{operation}. Ils proviennent uniquement du relevé statistique saisi en
E11. Il ne faut donc pas les comparer aux 20 opérations de la partie C : ce
sont deux jeux de données indépendants, destinés à des démonstrations
différentes.
\end{remarquebox}

\section{E12 : vue objet \code{Agence\_VOR}}

La vue objet \code{Agence\_VOR} expose une agence avec deux collections : les
conseillers actuellement affectés et les comptes courants de l'agence. Elle
possède un identifiant objet fondé sur le code de l'agence, ce qui permet de
l'utiliser comme un objet dans les requêtes et dans la méthode
\code{mutation}.

\begin{lstlisting}[caption={Type de la vue objet \code{Agence\_VOR}.},label={lst:e12-type}]
CREATE TYPE agence_vor_type AS OBJECT (
  code            VARCHAR2(6),
  nom             VARCHAR2(40),
  ville           VARCHAR2(40),
  ref_responsable REF conseiller_type,
  conseillers     conseiller_nt_type,
  comptes         cptCourant_nt_type,
  MEMBER PROCEDURE mutation(
    nConseiller   NUMBER,
    nAgenceCible  VARCHAR2
  )
);
\end{lstlisting}

La vue est construite avec \code{CAST(MULTISET(...))}. La première
sous-requête sélectionne les conseillers dont l'affectation est ouverte et dont
l'agence correspond à \code{MAKE\_REF(agence, a.code)}. La seconde sélectionne
les comptes courants de la même agence. \code{MAKE\_REF} est possible parce que
la table \code{agence} a été déclarée avec
\code{OBJECT IDENTIFIER IS PRIMARY KEY} : Oracle peut fabriquer une
\code{REF} à partir de la clé.

\begin{lstlisting}[caption={Création de la vue objet \code{Agence\_VOR}.},label={lst:e12-vue}]
CREATE OR REPLACE VIEW Agence_VOR OF agence_vor_type
WITH OBJECT IDENTIFIER (code) AS
SELECT a.code, a.nom, a.ville, a.ref_responsable,
       CAST(MULTISET(
         SELECT VALUE(c)
           FROM conseiller c, affectation af
          WHERE af.ref_conseiller = REF(c)
            AND af.ref_agence = MAKE_REF(agence, a.code)
            AND af.date_fin IS NULL
       ) AS conseiller_nt_type) AS conseillers,
       CAST(MULTISET(
         SELECT VALUE(cc)
           FROM cpt_courant cc
          WHERE cc.ref_agence = MAKE_REF(agence, a.code)
       ) AS cptCourant_nt_type) AS comptes
  FROM agence a;
\end{lstlisting}

Les collections se lisent avec \code{TABLE(v.conseillers)} et
\code{TABLE(v.comptes)}. La fonction \code{CARDINALITY} fonctionne ici, car il
s'agit de \emph{nested tables} et non de \code{VARRAY}s. Elle permet de
compter les conseillers et les comptes par agence sans déplier les collections.

\begin{lstlisting}[caption={Interrogation des collections de la vue objet.},label={lst:e12-requetes}]
SELECT v.code, k.nom, k.prenom
  FROM Agence_VOR v, TABLE(v.conseillers) k;

SELECT v.code, c.num_compte, c.solde
  FROM Agence_VOR v, TABLE(v.comptes) c;

SELECT v.code, v.nom,
       CARDINALITY(v.conseillers) AS nb_conseillers,
       CARDINALITY(v.comptes)     AS nb_comptes
  FROM Agence_VOR v
 ORDER BY v.code;
\end{lstlisting}

\section{La méthode \code{mutation}}

La méthode \code{mutation} reçoit un numéro de conseiller et un code d'agence
cible. Elle vérifie que le conseiller existe et qu'il est actuellement affecté
à l'agence source. Si ce n'est pas le cas, elle lève \code{ORA-20010}. Si le
conseiller ou l'agence cible est introuvable, elle lève \code{ORA-20011}. Elle
signale ensuite si le conseiller est responsable de l'agence qu'il quitte,
ferme l'affectation courante et en ouvre une nouvelle.

\begin{lstlisting}[caption={Méthode \code{mutation} de la vue objet.},label={lst:e12-mutation}]
MEMBER PROCEDURE mutation(nConseiller NUMBER, nAgenceCible VARCHAR2) IS
  v_ref_cons    REF conseiller_type;
  v_ref_ag_cur  REF agence_type;
  v_ref_ag_cib  REF agence_type;
  v_resp        NUMBER := 0;
  v_nb          NUMBER := 0;
BEGIN
  SELECT REF(c) INTO v_ref_cons
    FROM conseiller c WHERE c.num_conseiller = nConseiller;

  SELECT REF(a) INTO v_ref_ag_cur
    FROM agence a WHERE a.code = SELF.code;

  SELECT REF(a) INTO v_ref_ag_cib
    FROM agence a WHERE a.code = nAgenceCible;

  SELECT COUNT(*) INTO v_nb
    FROM affectation af
   WHERE af.ref_conseiller = v_ref_cons
     AND af.ref_agence = v_ref_ag_cur
     AND af.date_fin IS NULL;

  IF v_nb = 0 THEN
    RAISE_APPLICATION_ERROR(-20010,
      'Le conseiller n''est pas en poste dans l''agence source.');
  END IF;

  SELECT COUNT(*) INTO v_resp
    FROM agence a
   WHERE a.code = SELF.code
     AND a.ref_responsable = v_ref_cons;

  IF v_resp > 0 THEN
    DBMS_OUTPUT.PUT_LINE('ALERTE : le conseiller est responsable '
                         ||'de l''agence qu''il quitte.');
  END IF;

  UPDATE affectation
     SET date_fin = GREATEST(SYSDATE, date_debut) + 1/86400
   WHERE ref_conseiller = v_ref_cons
     AND ref_agence = v_ref_ag_cur
     AND date_fin IS NULL;

  INSERT INTO affectation (id_affectation, ref_conseiller, ref_agence,
                           date_debut, date_fin)
  VALUES (seq_affectation.NEXTVAL, v_ref_cons, v_ref_ag_cib, SYSDATE, NULL);
EXCEPTION
  WHEN NO_DATA_FOUND THEN
    RAISE_APPLICATION_ERROR(-20011, 'Conseiller ou agence introuvable.');
END;
\end{lstlisting}

\begin{attentionbox}[Pourquoi ajouter une seconde à \code{date\_fin} ?]
La contrainte \code{CK\_AFF\_DATES} impose \code{date\_fin > date\_debut}. Si
la mutation a lieu le jour même de l'ouverture de l'affectation,
\code{SYSDATE} peut être égal à \code{date\_debut}, ce qui viole la contrainte
et produit \code{ORA-02290}. L'expression
\code{GREATEST(SYSDATE, date\_debut) + 1/86400} garantit une différence d'au
moins une seconde. C'est une correction simple, mais indispensable pour les
tests de mutation immédiate.
\end{attentionbox}

\begin{clebox}[Le responsable n'est pas remplacé]
La méthode signale que le conseiller quitte une agence dont il est responsable,
mais elle ne modifie pas \code{agence.ref\_responsable}. Le sujet demande de
signaler le changement, pas de changer le responsable. Cette décision relève
d'une règle de management qui ne peut pas être automatisée sans risque. La
méthode laisse donc la décision au gestionnaire.
\end{clebox}

\section{Résultats E12}

La vue contient trois lignes : A1, A2 et A3. A1 a un conseiller et deux comptes
(CC4, CC7) ; A2 a deux conseillers et deux comptes (CC2, CC5) ; A3 a un
conseiller et deux comptes (CC3, CC6). Les mutations réussies déplacent KACI
(4) de A2 vers A3, puis BENALI (1) de A1 vers A2. Les tests de restauration
ramènent 1 vers A1 et 4 vers A2. Les compteurs de conseillers sont donc
restaurés : A1 = 1, A2 = 2, A3 = 1.

\begin{table}[htbp]
  \centering
  \small
  \caption{Résultats des tests de la vue objet et de la mutation.}
  \label{tab:resultats-e12}
  \begin{tabularx}{\textwidth}{@{}L{6.4cm}L{4.5cm}X@{}}
    \toprule
    \entete{Test} & \entete{Résultat attendu} & \entete{Commentaire} \\
    \midrule
    Vue \code{Agence\_VOR} & 3 lignes : A1, A2, A3 &
    Une ligne par agence. \\
    Conseillers A1/A2/A3 & 1 / 2 / 1 & Affectations ouvertes. \\
    Comptes A1/A2/A3 & 2 / 2 / 2 & CC4 et CC7 ; CC2 et CC5 ; CC3 et CC6. \\
    KACI (4) de A2 vers A3 & Succès & Mutation d'un non-responsable. \\
    BENALI (1) de A1 vers A2 & Succès, avec alerte &
    Le conseiller est responsable de A1. \\
    \code{A3.mutation(1, 'A2')} & \code{ORA-20010} &
    Le conseiller 1 n'est pas affecté à A3. \\
    \code{mutation(99, 'A2')} & \code{ORA-20011} &
    Conseiller inexistant. \\
    \code{mutation(1, 'A9')} & \code{ORA-20011} &
    Agence cible inexistante. \\
    \bottomrule
  \end{tabularx}
\end{table}

Après les tests, la table \code{affectation} contient neuf lignes : cinq
affectations initiales et quatre mutations historisées. Quatre lignes sont
fermées, cinq sont ouvertes. La vue ne compte que les affectations ouvertes.
L'historique est donc conservé, et la vue reste cohérente avec l'état courant.

\begin{attentionbox}[Rejouer E12]
La vue et le type \code{agence\_vor\_type} sont supprimés en début de partie E.
Lorsqu'on réutilise le même identifiant objet dans une nouvelle transaction, il
faut valider la transaction précédente par \code{COMMIT}. Sinon, Oracle peut
conserver l'ancien OID et refuser la réutilisation. Les scripts de test valident
après chaque mutation pour éviter ce problème.
\end{attentionbox}

\section{Pièges de la partie E}

\begin{table}[htbp]
  \centering
  \small
  \caption{Pièges de la partie E et solutions.}
  \label{tab:pieges-E}
  \begin{tabularx}{\textwidth}{@{}L{5.1cm}L{4.3cm}X@{}}
    \toprule
    \entete{Piège} & \entete{Symptôme} & \entete{Solution} \\
    \midrule
    \code{CARDINALITY} sur \code{mois\_vry} & \code{ORA-00932} &
    \code{TABLE(...)} + \code{COUNT(*)}. \\
    \code{TABLE(m.\allowbreak COLUMN\_VALUE.\allowbreak types)} &
    \code{ORA-00904} &
    \code{TABLE(m.types)}. \\
    \code{date\_fin = SYSDATE} & \code{ORA-02290} sur
    \code{CK\_AFF\_DATES} &
    \code{GREATEST(SYSDATE, date\_debut) + 1/86400}. \\
    Mutation sur un conseiller absent & \code{ORA-20010} &
    C'est un test attendu. \\
    Mutation sur un inconnu & \code{ORA-20011} &
    C'est un test attendu. \\
    \code{MAKE\_REF} sans OID basé sur la PK &
    \code{ORA-22979} ou impossible &
    \code{agence} a \code{OBJECT IDENTIFIER IS PRIMARY KEY}. \\
    Réutilisation d'un OID sans \code{COMMIT} &
    Ancien OID conservé &
    Valider la transaction avant de réutiliser l'identifiant. \\
    \bottomrule
  \end{tabularx}
\end{table}

\section{Conclusion de la partie E}

La partie E montre les mécanismes objet les plus expressifs du projet. Le type
récursif permet de naviguer dans une hiérarchie de parrainage sans table
supplémentaire. La collection multi-niveaux stocke un relevé annuel complet
dans une seule colonne, tout en restant interrogeable par double
\code{TABLE}. La vue objet agrège des collections calculées et expose une
méthode qui modifie l'historique des affectations. Ces trois démonstrations
complètent l'héritage, les \code{REF} et les méthodes des parties précédentes.
